% ═══════════════════════════════════════════════════════════
% §2 Results
% ═══════════════════════════════════════════════════════════
\section{Results}
\label{sec:results}

\subsection{Universal contact circuits transfer across proteins}
\label{sec:res-transfer}

We trained Quine--McCluskey-minimised Boolean circuits on consensus-residue
identities and separation distances from approximately 100 proteins drawn
from the PSICOV150 benchmark~\cite{jones2012}, then evaluated their ability
to predict native contacts on 50 entirely held-out proteins sharing no
evolutionary history with the training set (maximum pairwise 5-mer Jaccard
index 0.029 across all $\binom{150}{2}$ pairs).  The identity-level circuit
achieved a mean top-$L/5$ precision of \textbf{0.179} on held-out proteins,
representing $5.1\times$ enrichment over the random expectation of 0.035 and
a $1.9\times$ improvement over plain mutual information computed per-protein
(0.097).  The group-level circuit achieved 0.142.

At the circuit-membership threshold, the binary classifier achieved MCC
$\approx 0.095$ with enrichment of $2.5$--$2.8\times$.  Sensitivity analysis
confirmed that ranked precision was invariant across all nine combinations of
labelling thresholds tested ($T \in \{1.5, 2, 3\} \times n_{\min} \in \{10,
	30, 60\}$).

The minimised circuits produced biophysically interpretable rules.  Essential
prime implicants included hydrophobic$\times$aromatic packing at separations
6--21 residues, sulfur-mediated contacts at separations 6--11, and specific
Level-2 vocabularies such as $\{M,F\}\times\{M,F,Y,W\}$ consistent with
methionine--phenylalanine core interactions.  These rules were rediscovered
from sequence data without any biophysical prior knowledge.

For context, the published PSICOV predictions achieve a top-$L/5$ precision
of 0.727 on the same held-out proteins when evaluated against the same native
contact labels, confirming both the quality of our ground-truth pipeline
($\Delta = 0.003$ from Jones et al.'s Table~1) and the magnitude of the gap
between interpretable discrete circuits and continuous sparse-inverse-covariance
methods.

\subsection{Gray-code adjacency is enriched among native contacts}
\label{sec:res-gray}

The Karnaugh-map framework predicts that residue pairs connected by a single
bit-flip in Gray-code space should be enriched among native contacts, because
the encoding maps chemically similar amino acids to adjacent codewords.
Among 47{,}430 native-contact pairs across 150 PSICOV proteins, we found
that 22.83\% are Gray-adjacent ($d_G = 1$), compared to a
composition-controlled background rate of 19.37\%, yielding an enrichment of
$1.178\times$ with a two-sided exact binomial $p = 1.1 \times 10^{-77}$.
Bootstrap confidence intervals from protein-level resampling (2{,}000
replicates) span [22.27\%, 23.40\%].

A permutation test over random relabelings of the twenty amino-acid codes
demonstrated that this enrichment is specific to the He/Gray assignment:
random code permutations yield a mean enrichment of $0.961 \pm 0.095$, placing
the true encoding at $z = 2.51$ above the null distribution.  This confirms
that the particular Gray-code structure carries genuine signal beyond what
generic physicochemical clustering would produce, although most of the raw
enrichment reflects the broader tendency of chemically similar residues to
contact each other.

An earlier analysis restricted to the SARS-CoV-2 Spike protein found an
identical point estimate ($1.18\times$) but failed to reach significance
($p = 0.345$) because only 38 contact pairs were available.  The present
analysis with three orders of magnitude more data resolves this ambiguity:
the effect is real but requires adequate statistical power to detect.

\subsection{Rank fusion improves DCA ranking}
\label{sec:res-fusion}

The Karnaugh-map cell-rate prior improved global ranking (circuit AUC 0.658
vs.\ DCA AUC 0.583) while degrading top-$K$ sharpness relative to raw DCA.
We hypothesised that fusing the two signals via rank aggregation would capture
both properties.  Five fusion rules were evaluated under identical protein-level
cross-validation:

\begin{table}[ht]
	\centering
	\caption{Rank-fusion results over five-fold protein-level cross-validation.}
	\label{tab:fusion}
	\begin{tabular}{lcc}
		\toprule
		Method            & prec@L/5 (mean $\pm$ sd)   & AUC            \\
		\midrule
		rawDCA alone      & $0.169 \pm 0.125$          & 0.583          \\
		circH prior alone & $0.124 \pm 0.088$          & \textbf{0.658} \\
		\midrule
		\textbf{ranksum}  & $\mathbf{0.204 \pm 0.104}$ & 0.638          \\
		geometric mean    & $0.203 \pm 0.103$          & 0.634          \\
		product           & $0.199 \pm 0.108$          & 0.592          \\
		\bottomrule
	\end{tabular}
\end{table}

The rank-sum fusion achieved prec@L/5 = \textbf{0.204}, exceeding rawDCA
(Wilcoxon $p < 10^{-4}$, mean difference $+0.036$) and circH ($p < 10^{-4}$,
mean difference $+0.080$).  Cohen's $d = 0.31$ with a bootstrap 95\% CI of
$[0.017, 0.054]$ on the paired difference.  Mechanistic decomposition revealed
that all rescued pairs had above-median scores from \emph{both} methods,
confirming coincidence-of-evidence amplification rather than rescue of weak
signals.  Adding plain MI as a third signal did not improve further
(three-way 0.201 vs.\ two-way 0.204, $p = 0.607$).

This constitutes the first demonstration that Karnaugh-map-derived features
carry structural information orthogonal to continuous coupling scores from
DCA, providing a novel post-processing step applicable to any existing
coevolution pipeline.

\subsection{Flipped circuits generalise across proteins}
\label{sec:res-flipped}

In addition to positive (mutation-detecting) maps, we constructed flipped maps
in which ON cells represent never-observed residue combinations, i.e. negative
selection constraints.  When trained on PSICOV proteins and applied to
held-out proteins, the flipped circuit achieved a specificity of
\textbf{98.3\%}: forbidden-called cells showed exactly the predicted
$\leq 1.75\%$ contact rate on unseen proteins, matching the training-set
upper bound almost perfectly.

Combined with the positive channel, the two-sided classifier achieved MCC =
$+0.169$ over definitive calls (22\% coverage), confirming that both polarity
channels carry complementary structural information.

\subsection{Conserved column pairs are contact-enriched}
\label{sec:res-conservation}

Pairs in which both alignment columns are fully conserved (entropy $< 0.001$)
showed a native-contact rate of 11.4\%, representing a $3.2\times$ enrichment
over the overall base rate of 3.5\%.  However, a finer analysis using
continuous minimum-pair entropy revealed a non-monotone gradient: absolutely
invariant positions (entropy $\approx 0$) were depleted of contacts
($0.69\times$), while near-conserved positions (entropy $0.001$--$0.05$) were
enriched ($1.35\times$).

This non-monotonicity is biologically interpretable: absolutely invariant
positions are frequently catalytic residues or post-translational modification
sites whose conservation reflects functional constraint rather than structural
packing.  Near-conserved positions are sufficiently constrained to maintain
structural integrity while retaining enough variability to generate detectable
covariation signal.  The overall Spearman correlation between column entropy
and contact status is negative ($\rho = -0.053$, $p \approx 0$), confirming
the monotone trend.

\subsection{Contact maps are more Karnaugh-map-compressible than chance}
\label{sec:res-compressibility}

For six PSICOV proteins small enough for exact Quine--McCluskey minimisation
of a padded $64 \times 64$ grid, we compared the compression ratio (\# prime
implicants / \# contact cells) between real contact matrices and controls whose
labels were shuffled within separation bins.  Real maps required fewer prime
implicants than shuffled controls in all six cases (mean 2.31 vs 2.74;
Wilcoxon $p = 0.031$), demonstrating that three-dimensional geometry carries
Karnaugh-map-expressible regularity beyond what density and separation
structure alone explain.

\subsection{Reconstruction assessment}
\label{sec:res-reconstruction}

Against literature thresholds for structure reconstruction readiness (median
long-range precision@L/5 $\geq 0.30$ AND recall@L $\geq 0.10$, derived from
Jones et al.'s observation that $\geq 0.5$ long-range precision suffices for
fold determination~\cite{jones2012}), our circuits do not meet the criterion
(median LR prec@L/5 = 0.000; recall@L = 0.028).  Published PSICOV predictions
do meet it (LR prec@L/5 = 0.683; recall@L = 0.218).  At current signal
strength, the circuits serve as an annotation layer identifying specific
coevolutionary constraints; they do not supply sufficient spatial information
for template-free structure determination.

\subsection{Negative results constrain the theory}
\label{sec:res-negative}

Three hypotheses were tested and rejected.  First, an attention-style tension
mechanism normalising each position's mutual-information profile by its
strongest partner did not improve circuit precision (tension-augmented 0.154
vs.\ non-tension 0.165), and high-tension positions were anti-correlated with
structural hubs (median $\rho = -0.06$).  Second, quartile-binned DCA scores
added as Karnaugh-map feature dimensions provided no improvement over identity
features alone (circH 0.124 vs L2 0.165), indicating that continuous coupling
information cannot be productively discretised into the Boolean framework.
Third, iterative direct information from our mean-field implementation was
anti-predictive on families with many conserved columns due to pseudo-inverse
conditioning, leading us to adopt APC-corrected Frobenius norm as the primary
DCA score.

Each negative result constrains where the framework applies and directs future
effort toward productive extensions rather than unproductive pursuit of
approaches that the data do not support.
